\chapter{Topological Nucleation, the Fossil CMB, and Lattice Annealing}
\label{chap:nucleation_annealing}

The preceding chapter established Catastrophic Decompression as the mechanical rupture that restores the active State-I FCC vacuum from the globally percolated State-II configuration. The present chapter follows the immediate physical consequences of that event.

Standard cosmology describes the early universe as a hot, dense particle--radiation plasma, with inflation commonly invoked in standard early-universe models to account for large-scale homogeneity and recombination establishing the photon-decoupling and last-scattering epoch associated with the Cosmic Microwave Background (CMB). Tension Field Theory (TFT) proposes a fundamentally different mechanical chronology.

TFT proposes that the rupture seeds three coupled but distinct outcomes: localized topological matter defects, a chiral nucleation bias between matter and antimatter, and a highly excited radiation--lattice state that subsequently anneals toward transparency. The thermal radiation surviving that early evolution is then treated as the primordial boundary of the Fossil-CMB problem.

The chronology is therefore not a set of simultaneous processes, but a mechanical sequence:
\begin{equation}
\begin{aligned}
    \text{Catastrophic Decompression} 
    &\longrightarrow \text{topological nucleation} 
    \longrightarrow \text{lattice--radiation annealing} \\
    &\longrightarrow \text{transparency} 
    \longrightarrow \text{Fossil CMB evolution}.
\end{aligned}
\end{equation}

The microscopic chiral branch is developed in \hyperref[app:baryon_asymmetry]{Appendix~\ref*{app:baryon_asymmetry}}, the conditional annealing timescale is developed in \hyperref[app:lattice_annealing]{Appendix~\ref*{app:lattice_annealing}}, and the post-transparency thermodynamic requirements of the Fossil-CMB route are developed in \hyperref[app:cmb_thermal_completion]{Appendix~\ref*{app:cmb_thermal_completion}}.


\section{Topological Nucleation (The Birth of Matter)}

In TFT, matter does not enter the active vacuum as an independent material substance. As established in \hyperref[def:definition_three]{Part I}, a massive particle is a persistent localized topological defect of the State-I FCC vacuum.

Its inertial loading is represented microscopically by
\begin{equation}
m_{\mathrm{i}} = M_{\mathrm{top}} = \rho_{\mathrm{IC}}^{(\mathrm{imp})} V_{\mathrm{imp}} f_{\varphi}(\Delta\varphi),
\end{equation}
where the defect--condensate impedance construction is developed in \hyperref[app:condensate_impedance]{Appendix~\ref*{app:condensate_impedance}}.

Catastrophic Decompression supplies a natural environment in which such nontrivial defects may be nucleated. The rupture violently converts the globally saturated phase-boundary network back into the active FCC configuration. During this rapid reconfiguration, not every localized phase distortion need relax into the topologically trivial vacuum state.

TFT therefore proposes that some rupture-generated boundary distortions close into protected localized phase configurations and survive as primordial matter defects. In this mechanical sense, matter is the topological ``shrapnel'' of Catastrophic Decompression: not fragments of a separate substance, but stable nontrivial configurations left behind as the active vacuum reforms.

This interpretation is stronger than a purely energetic creation picture, but its status must remain precise. The existence and inertial loading of localized topological matter states are developed independently in Parts I and II; the complete dynamical calculation showing how Catastrophic Decompression produces the observed proton and neutron population has not yet been derived. The rupture-to-baryon nucleation mechanism is therefore a physical TFT hypothesis constrained by the established defect mechanics, not a completed particle-production calculation.

The global baryonic energy allocation is likewise not derived here, but it is strictly bounded by the lattice geometry. In the present TFT cosmological ledger, an approximate baryonic structural fraction is assigned to localized defects:
\begin{equation}
    \beta_{\mathrm b}
    \approx
    0.05
\end{equation}
Independently, as developed in \hyperref[app:cosmic_energy_budget_and_chronology]{Appendix~\ref*{app:cosmic_energy_budget_and_chronology}}, the FCC Octahedral-site construction yields an upper isotropic geometric accommodation limit:
\begin{equation}
    \beta_{\mathrm{oct}}
    =
    \frac{\pi(\sqrt{2}-1)^3}{3\sqrt{2}}
    \approx
    0.0526
\end{equation}

The numerical proximity of these quantities is a vital structural consistency relation: the Octahedral limit provides an upper isotropic geometric accommodation scale for condensed defect states. It does not, by itself, derive the observed baryonic fraction or prove that every baryonic defect is uniquely confined to a single microscopic Octahedral site.

The spatial distribution of the primordial defects is a separate question. Because the Big Rupture is a distributed failure of a reservoir-scale phase-boundary network rather than an explosion from a single material point, TFT admits distributed nucleation throughout the newly restored active domain. Whether this mechanism quantitatively reproduces the observed large-scale matter distribution remains an observational requirement rather than an assumed consequence of the rupture.


\vspace{0.5em}

\noindent\textbf{Instantaneous Core Genesis (\(t \approx 0^{+}\) Primordial State-II Seeds)}

Primordial matter defects and primordial State-II cores obey different formation conditions. Atomic binding requires the post-rupture medium to anneal toward the neutral-hydrogen stability regime, whereas State-II formation is controlled by local structural saturation.

Immediately after Catastrophic Decompression, strong compression and interference among rupture-generated stress transients can locally drive the restored FCC vacuum back toward the convergence boundary:
\begin{equation}
    \rho
    \rightarrow
    \rho_{\mathrm{QG}},
    \qquad
    \alpha_{\mathrm G}
    \rightarrow
    1.
\end{equation}

At this boundary the active FCC reconfiguration capacity vanishes and the local configuration phase-locks into State II. TFT therefore permits prompt post-rupture State-II nucleation, providing a candidate origin for the primordial seed population required by the later macroscopic cycle.

The present section does not determine the number, mass spectrum, or spatial distribution of those primordial seeds. Their aggregate initial condition and subsequent role in the global structural clock are deferred to \hyperref[ch:macroscopic_oscillator]{Chapter~\ref*{ch:macroscopic_oscillator}} and \hyperref[app:cosmic_energy_budget_and_chronology]{Appendix~\ref*{app:cosmic_energy_budget_and_chronology}}.

The thermal chronology is independent of this saturation criterion. Under the conditional annealing closure developed in \hyperref[app:lattice_annealing]{Appendix~\ref*{app:lattice_annealing}}, the post-rupture system reaches the adopted \(T_{\mathrm a}\sim3000~\mathrm K\) neutral-hydrogen stability and transparency boundary after approximately
\begin{equation}
    t_{\mathrm a}
    \simeq
    3.86\times10^{5}\ \mathrm{yr}.
\end{equation}

Thus, primordial State-II seeds may form essentially at the rupture boundary, while stable neutral atomic matter belongs to the later annealed epoch.

\section{Chiral Nucleation and Baryon Asymmetry}

The observed dominance of matter over antimatter remains one of the major unresolved problems of fundamental cosmology. Although the Standard Model contains the ingredients required for baryon-number-generating processes, no experimentally established mechanism yet accounts for the observed cosmic asymmetry. Tension Field Theory (TFT) proposes a mechanical origin by applying angular-momentum conservation to the topological phase transition of the spacetime lattice.

The universe did not unfold from a formless, isotropic explosion. As established in the preceding chapter, the Big Rupture originated from a globally percolated network of State-II domains. If this pre-rupture macroscopic fault network carried a coherent net angular momentum, TFT proposes that part of this conserved angular momentum was inherited by the newly activated Soft Floor as a transient torsional field---a lattice twist---during decompression. The emerging State-I FCC lattice would then possess a temporary preferred helical ``grain'' during the nucleation epoch.

Within TFT, matter and antimatter are represented as opposite chiral topological knots. In a torsionally neutral vacuum, their formation barriers are taken to be degenerate at the base energy \(E_{\mathrm{base}}\). A transient torsional strain lifts this degeneracy:

\begin{itemize}
    \item \textbf{Matter (aligned with the torsion):} Forms with the helical grain of the active vacuum. The lattice twist assists formation of the aligned topological knot, giving the lower barrier
    \begin{equation}
        E_{\mathrm{matter}} = E_{\mathrm{base}} - E_{\mathrm{torsion}}.
    \end{equation}

    \item \textbf{Antimatter (anti-aligned):} Forms against the helical grain. The lattice resists formation of the oppositely oriented knot, giving the higher barrier
    \begin{equation}
        E_{\mathrm{antimatter}} = E_{\mathrm{base}} + E_{\mathrm{torsion}}.
    \end{equation}
\end{itemize}

Under the quasi-thermal nucleation approximation developed in \hyperref[app:baryon_asymmetry]{Appendix~\ref*{app:baryon_asymmetry}}, the relative population of the two chiral sectors is then governed by the two-state Boltzmann ratio,
\begin{equation}
    \frac{N_{\mathrm{matter}}}{N_{\mathrm{antimatter}}} = \exp\left( \frac{2E_{\mathrm{torsion}}}{k_{\mathrm B}T_{\mathrm n}} \right).
\end{equation}

Equivalently, defining the reduced matter--antimatter number asymmetry,
\begin{equation}
    \mathcal A_{\mathrm B} \equiv \frac{N_{\mathrm{matter}} - N_{\mathrm{antimatter}}}{N_{\mathrm{matter}} + N_{\mathrm{antimatter}}},
\end{equation}
gives
\begin{equation}
    \mathcal A_{\mathrm B} = \tanh\left( \frac{E_{\mathrm{torsion}}}{k_{\mathrm B}T_{\mathrm n}} \right) \simeq \frac{E_{\mathrm{torsion}}}{k_{\mathrm B}T_{\mathrm n}}
\end{equation}
for
\begin{equation}
    E_{\mathrm{torsion}} \ll k_{\mathrm B}T_{\mathrm n}.
\end{equation}

The result exposes the central mechanical requirement directly: a primordial reduced chiral asymmetry of order \(10^{-9}\) requires only a torsional formation-energy splitting of order
\begin{equation}
    \frac{E_{\mathrm{torsion}}}{k_{\mathrm B}T_{\mathrm n}} \sim 10^{-9}.
\end{equation}

Thus an extremely small structural bias can, in principle, be amplified into a surviving matter excess during the nucleation and annihilation epoch. Antimatter is not required to be absent from the primordial state; in this mechanism it is statistically disadvantaged during crystallization of the active vacuum. The absolute value of the torsional bias and its conversion into the observed baryon-to-photon abundance remain quantities to be derived from the complete nucleation and thermal history.

\vspace{0.5em}

\noindent\textbf{Residual Torsion and the Cosmic-Birefringence Constraint}
\label{sec:resolution_cosmic_birefringence}

A primordial torsional bias strong enough to separate the two chiral nucleation sectors cannot persist unchanged into the post-annealing universe. Any residual parity-odd optical coupling is tightly constrained by CMB polarization measurements of cosmic birefringence \cite{Minami_2020}.

TFT therefore proposes that the rupture-generated torsional loading is substantially transferred into the newly formed chiral defect population during nucleation. In this picture, the enormous population of topological matter defects acts as a ``topological sponge,'' converting part of the transient macroscopic lattice twist into localized chiral structure.

This transfer mechanism is physically natural within the TFT topology, but its quantitative angular-momentum and torsion ledger has not yet been derived. Cosmic birefringence therefore provides a direct downstream falsification constraint: a successful microscopic completion must show that the residual post-nucleation torsional field relaxes to a level consistent with CMB polarization bounds while retaining the chiral bias required by the matter--antimatter asymmetry.

\section{The Fossil Echo: The Mechanical Origin of the CMB}
\label{sec:phonon_echo_cmb}

Within Tension Field Theory (TFT), the Cosmic Microwave Background is a thermal electromagnetic population occupying the active photon branch. Its primordial boundary is identified, within the present TFT closure, with the Catastrophic Decompression and subsequent lattice-annealing history. For the isotropic thermal photon field, the energy density is given by
\begin{equation}
    u_{\mathrm{CMB}}(T) = \frac{\pi^2 k_{\mathrm B}^4}{15\hbar^3 c_{\mathrm E}^3} T^4.
\end{equation}
The microscopic thermal closure of this population, distinguishing it from nonradiative Active-Vacuum receiving modes ($u_{\mathrm{CMB}} \neq u_{\mathrm{AV}}$), is developed in \hyperref[app:cmb_thermal_completion]{Appendix~\ref*{app:cmb_thermal_completion}}.

Immediately following the Big Rupture, Catastrophic Decompression violently excites the newly restored State-I medium. Nonlinear pressure fronts, shear, phase-boundary recoil, and compressive acoustic shockwaves propagate through the forming Active Vacuum while newly nucleated baryonic shrapnel is redistributed by this far-from-equilibrium environment. The rupture-generated acoustic structure provides the mechanical mechanism for imprinting large-scale density correlations in matter.

While the macroscopic matter distribution is locked in by these acoustic fronts, the early radiation--matter system undergoes the separate annealing history developed in \hyperref[app:lattice_annealing]{Appendix~\ref*{app:lattice_annealing}}. Under the diffusive bulk-cooling similarity and geometric annealing-mode closure, the system reaches the adopted neutral-hydrogen stability and transparency scale:
\begin{equation}
    T_{\mathrm a} \sim 3000\ \mathrm K, \qquad t_{\mathrm a} \simeq 3.86\times10^5\ \mathrm{yr}.
\end{equation}

The thermal photon population present at this boundary defines the earliest currently defensible Fossil-CMB state in TFT. Its subsequent evolution requires energy exchange with independently specified finite-energy receiving modes together with photon-number-changing detailed balance to maintain an exact Planck spectrum.

The primordial chronology therefore proceeds as a strict mechanical sequence:

\begin{equation}
\begin{aligned}
    \text{Catastrophic Decompression}
    &\longrightarrow
    \begin{cases}
        \text{mechanical density imprint},\\
        \text{primordial thermal radiation},
    \end{cases}
    \\
    &\longrightarrow
    \text{annealing and transparency}
    \\
    &\longrightarrow
    \text{Fossil CMB}.
\end{aligned}
\end{equation}

The surviving Fossil CMB establishes the ambient electromagnetic baseline of the present Active Vacuum. We define this radiative baseline as the \textbf{CMB Noise Floor}: the statistically isotropic thermal photon background against which a progressively relaxed coherent photon excitation ultimately loses its distinct observational identity. Finite-energy receiving modes of the Active Vacuum mediate the transfer of coherent carrier energy, while the CMB Noise Floor represents the radiative population into which that transferred energy can participate through thermal absorption--emission exchange.

The Photon Energy-Relaxation Length $R_{\mathrm U}$ sets the characteristic one-$e$-fold relaxation length of the coherent photon branch. Continued propagation beyond $R_{\mathrm U}$ drives the surviving carrier progressively closer to the CMB Noise Floor as its coherent energy is transferred into the surrounding Active-Vacuum environment.

The present $2.725$ K background is consequently treated as the surviving thermal photon field of this primordial history and as the present electromagnetic realization of the CMB Noise Floor. Its Planckian character can be obtained from a genuine number-changing thermal interaction, while the microscopic thermal closure connecting the receiving sector and the Fossil photon population is developed in
\hyperref[app:cmb_thermal_completion]{Appendix~\ref*{app:cmb_thermal_completion}}.

\subsection{The Cosmic Web: Ghost Voids and Superclusters}
\label{sec:cosmic_web_formation}

The macroscopic acoustic shells generated by the Big Rupture do not simply fade away; their matter-displacement pattern becomes the fossil structural template of the observable Cosmic Web.

Before Catastrophic Decompression, inertial matter associated with each State II domain is phase-segregated onto its live FCC/SC boundary, while the saturated SC interior itself contains no ordinary inertial matter. When global confinement fails at the Big Rupture, the compressed State II domains undergo rapid SC$\rightarrow$FCC decompression, driving their boundary-segregated baryonic matter outward with the expanding phase fronts.

\begin{itemize}

    \item \textbf{Ghost Voids (Acoustic Evacuation Zones):}
    
    Large cosmic voids are interpreted as fossilized interiors of decompressed State II domains. As the SC core reopens into the dynamic State I FCC vacuum, the baryonic matter formerly concentrated on its live boundary is swept outward, leaving behind a restored Active-Vacuum region with strongly reduced topological-matter occupancy. The resulting void is therefore not empty space, but a matter-depleted region of the active vacuum tracing the former location and overlap geometry of a precursor State II domain.

    \item \textbf{Filaments and Superclusters (Shell-Intersection Structures):}
    
    Because the Big Rupture involves a distributed network of neighboring State II domains, many matter-bearing decompression shells expand simultaneously. Where neighboring shells intersect, outward-swept baryonic matter is compressed into shared boundaries, naturally producing wall-like and filamentary structures. Regions in which several shells converge accumulate substantially more matter and form dense junctions that provide the mechanical precursors of supercluster-scale structures, including systems such as the Laniakea Supercluster \cite{Tully_2014}.

\end{itemize}

The Big Rupture therefore imprints a macroscopic structural pattern into the primordial matter distribution: decompressed State II interiors become matter-depleted voids, while intersections of their outward-moving boundary matter generate filaments, walls, and dense supercluster nodes. Under the optical projection developed in the following subsection, this fossil structure is associated with the TFT Effective Acoustic Scale,
\begin{equation}
    r_{\mathrm{s}}
    \approx
    42.73\ \mathrm{Mpc}.
\end{equation}

In TFT, the observed Cosmic Web is therefore interpreted as a fossilized mechanical imprint of Catastrophic Decompression.

\subsection{Acoustic Scale and the Optical Projection}
\label{sec:acoustic_variance}

The rupture picture naturally suggests a characteristic spatial scale, but TFT must strictly distinguish an independently generated mechanical scale from one constructed through observational projection.

For a specified constant physical Vacuum Relaxation Rate $\Gamma_{\mathrm D}$, the characteristic coherent-relaxation length of the photon sector is

\begin{equation}
    R_{\mathrm U}
    =
    \frac{c_{\mathrm E}}{\Gamma_{\mathrm D}}.
\end{equation}

This is a one-$e$-fold coherent scale-relaxation length, not a physical wall, a universal photon stopping radius, or the radius of the Structural Reservoir.

For the Euclidean acoustic interpretation developed later in
\hyperref[chap:cosmological_anomalies]{Chapter~\ref*{chap:cosmological_anomalies}},
TFT introduces the conditional projection closure:

\begin{equation}
    D_{\mathrm{proj}}
    \equiv
    R_{\mathrm U}.
\end{equation}

Applying the observed CMB acoustic angular scale then gives the TFT Effective Acoustic Scale:

\begin{equation}
    r_{\mathrm s}^{(\mathrm{TFT})}
    =
    D_{\mathrm{proj}}\theta_*
    =
    R_{\mathrm U}\theta_*.
\end{equation}

Under the retained low-redshift reference specialization in which the observed Hubble-scale slope is relaxation dominated, the reference projection scale is
$R_{\mathrm U}\simeq4104.5\ \mathrm{Mpc}$.

Using

\begin{equation}
    R_{\mathrm U}
    \simeq
    4104.5\ \mathrm{Mpc},
    \qquad
    \theta_*
    \simeq
    0.01041,
\end{equation}

yields

\begin{equation}
    r_{\mathrm s}^{(\mathrm{TFT})}
    \simeq
    42.73\ \mathrm{Mpc}.
\end{equation}

This numerical scale is an important TFT phenomenological result, but its status is precise: it follows from the proposed $D_{\mathrm{proj}}=R_{\mathrm U}$ closure and the observed angular scale. It is not yet an independent prediction of the Catastrophic-Decompression dynamics. The corresponding Euclidean acoustic-projection closure is developed in
\hyperref[sec:euclidean_acoustic_resolution]{Section~\ref*{sec:euclidean_acoustic_resolution}}.

Within the rupture interpretation, the $42.73$ Mpc scale can therefore be treated as a candidate characteristic scale against which the predicted precursor-core distribution and fossil density pattern must ultimately be tested. A successful first-principles completion would derive an equivalent scale directly from the percolation, phase-conversion, and post-rupture mechanics rather than obtaining it only through the observed angular projection.

Observed voids and large-scale structures need not possess a single universal diameter. Variance can arise naturally from the distribution of precursor State-II domains, the overlap of neighboring rupture fronts, and subsequent gravitational evolution. TFT must derive this distribution rather than identify every structure with an identical spherical shell.

\vspace{0.5em}

\noindent\textbf{Global Homogeneity}

The Structural Reservoir extends vastly beyond the characteristic optical relaxation scale,

\begin{equation}
    R_{\mathrm U}
    \ll
    R_{\mathrm{Res}}.
\end{equation}

This distinction allows the observable domain to occupy only a small part of the complete TFT structural reservoir. Statistical homogeneity of the precursor network across that larger reservoir is a macroscopic closure used in
\hyperref[ch:macroscopic_oscillator]{Chapter~\ref*{ch:macroscopic_oscillator}}.
However, it does not imply that identical $42.73$ Mpc shells must repeat exactly throughout the full lattice.

Likewise, $R_{\mathrm U}$ is not a CMB-assimilation or transparency boundary. Coherent photons may propagate beyond one $R_{\mathrm U}$, and any terminal photon-branch condition must be obtained independently from the microscopic stability of that branch rather than from the CMB frequency or temperature.

\section{The Duality of Fossil Radiation and Active Thermal Equilibrium}
\label{sec:vacuum_radiation_duality}

To understand what an astronomical detector actually measures in TFT, we must distinguish between the coherent photon branch and the finite-energy degrees of freedom of the Active Vacuum with which that radiation exchanges energy.

Photons propagate through the State-I FCC vacuum as coherent electromagnetic eigenmodes. Under the scale-relaxation mechanism developed in \hyperref[app:photon_dispersion_relaxation]{Appendix~\ref*{app:photon_dispersion_relaxation}} and \hyperref[app:vacuumdamping_cosmicredshift]{Appendix~\ref*{app:vacuumdamping_cosmicredshift}}, when the coherent scale-relaxation channel is active, a fraction of their organized propagating energy is transferred into noncoherent finite-energy modes of the Active Vacuum. The local energy handoff is represented by

\begin{equation}
    \dot E_{\mathrm{rec}}
    =
    -\dot E_{\gamma}
    =
    \Gamma_{\mathrm D}E_{\gamma}.
\end{equation}

Here, $E_{\mathrm{rec}}$ denotes the energy accumulated in the finite-energy receiving sector.

The receiving sector consists of those finite-energy Active-Vacuum degrees of freedom that are dynamically accessible to the photon branch. These modes provide the distributed mechanical capacity through which radiative energy can be stored, redistributed, and exchanged thermally with the photon population.

The two sectors therefore perform complementary roles: coherent photons carry radiation across the vacuum, while the finite-energy Active-Vacuum modes provide the mechanical channel through which that radiation can participate in thermal equilibrium.

Therefore, an astronomical bolometer measuring the \(2.725\) K background is not, in the TFT interpretation, simply photographing a distant ``wall'' of ancient Big Bang light associated with the standard last-scattering surface. It is detecting the fossil electromagnetic photon population produced by the rupture-and-annealing history, while the Active Vacuum provides the microscopic medium through which that radiation propagates and thermally exchanges energy.

The mechanism by which this energy exchange is converted into and maintains a Planck blackbody population requires a separate photon-number-changing absorption--emission interaction satisfying microscopic detailed balance, developed in the following section, with the complete microscopic thermal closure subjected to strict falsification criteria in \hyperref[app:cmb_thermal_completion]{Appendix~\ref*{app:cmb_thermal_completion}}.

\section{Lattice Hysteresis and the Proposed Blackbody Equilibrium}

\label{sec:lattice_hysteresis}

When the irreversible photon--receiver channel is active, the energy transfer described above may be represented mechanically in TFT as lattice hysteresis. As a coherent photon eigenpacket propagates through the Active FCC vacuum, its transverse electromagnetic oscillation couples weakly to the dynamically accessible receiving sector and induces a cyclic deformation of the medium. As derived in \hyperref[app:photon_dispersion_relaxation]{Appendix~\ref*{app:photon_dispersion_relaxation}}, the irreversible component of this deformation produces a positive hysteretic work per optical cycle:

\begin{equation}
    \Delta E_{\mathrm{cycle}}
    =
    \int_{V_\gamma}
    \oint_{\mathrm{cycle}}
    \sigma_T\,d\varepsilon_T\,dV
    > 0.
\end{equation}

Here $V_\gamma$ is the instantaneous spatial support of the coherent photon mode, while $\sigma_T$ and $\varepsilon_T$ are its work-conjugate transverse stress and strain. In the conservative limit the stress--strain loop encloses no area and $\Delta E_{\mathrm{cycle}}=0$.

This hysteresis represents a small but cumulative transfer of organized photon energy into the receiving Active-Vacuum sector. Because the coupling is weak, the coherent photon branch can remain long-lived over cosmological distances while continuously exchanging energy with the medium that supports it.

The Active FCC vacuum contains the finite-energy receiving degrees of freedom through which radiative energy may be absorbed and redistributed when the corresponding photon--receiver coupling is active. Its contiguous, hyper-connected structure therefore provides the mechanical medium within which coherent radiation transports energy and noncoherent excitations supply the distributed thermal capacity required for equilibration. Thermal re-emission belongs to the corresponding photon-number-changing absorber--emitter interaction.

Over sufficiently long timescales, when the thermalized photon sector possesses a photon-number-changing absorption--emission interaction satisfying microscopic detailed balance,

\begin{equation}
    \frac{
        \Gamma_{\mathrm{em}}(\omega,T)
    }{
        \Gamma_{\mathrm{abs}}(\omega,T)
    }
    =
    \exp\left(
        -\frac{\hbar\omega}{k_{\mathrm B}T}
    \right),
\end{equation}

repeated absorption, redistribution, and re-emission drive the coupled photon--Active-Vacuum system toward detailed balance. TFT proposes that this process organizes the radiation field toward a Planck blackbody distribution and establishes an Active Thermal Equilibrium of the functioning State-I vacuum.

For a closed Fossil photon--receiver system, the common equilibrium temperature is governed by conservation of the combined radiative and receiving-sector energy:

\begin{equation}
    {
    u_{\mathrm{CMB}}(T_{\gamma,0})
    +
    u_{\mathrm{rec}}(T_{\mathrm{rec},0})
    =
    u_{\mathrm{CMB}}(T_{\mathrm{eq}})
    +
    u_{\mathrm{rec}}(T_{\mathrm{eq}}).
    }
\end{equation}

For specified receiving thermodynamics and initial loading, this relation determines the common equilibrium temperature $T_{\mathrm{eq}}$. A finite active-to-latent interface specialization developed in \hyperref[app:cmb_thermal_completion]{Appendix~\ref*{app:cmb_thermal_completion}} yields a conditional equilibrium endpoint near

\begin{equation}
    T_{\mathrm{eq}}
    \simeq
    2.726\ \mathrm{K}
\end{equation}

under its stated interface and initial-state conditions.

The physical picture is therefore not one of passive photon energy loss, but of continuous conversion between coherent radiative transport and the thermal degrees of freedom of the same underlying vacuum medium. The quantitative thermal-equilibrium closure, including the microscopic absorber--emitter interaction, the receiving-sector thermodynamics, and the conditions required for exact Planck equilibrium, is developed in \hyperref[app:cmb_thermal_completion]{Appendix~\ref*{app:cmb_thermal_completion}}.

\section{The Lattice Annealing Law (Deriving the 380,000-Year Epoch)}
\label{sec:lattice_annealing_law}

In TFT, the historical \(380{,}000\)-year epoch is interpreted fundamentally as the \textbf{thermal annealing timescale} of the Spacetime Gravitational Floor. Following Catastrophic Decompression, the enormous structural energy released by the snap-back of the medium is redistributed through the newly restored FCC Active Vacuum. The characteristic substrate transition scale,

\begin{equation}
    T_{\mathrm S}
    \approx
    1.416 \times 10^{32}\ \mathrm{K},
\end{equation}

sets the high-energy boundary of this structural annealing process. It represents the energy-equivalent temperature of the substrate transition rather than an ordinary photon blackbody temperature.

During the annealing regime developed in \hyperref[app:lattice_annealing]{Appendix~\ref*{app:lattice_annealing}}, the diffusive bulk-cooling similarity gives

\begin{equation}
    T \propto t^{-1/2}.
\end{equation}

This annealing similarity applies to the pre-transparency structural-cooling regime and is not continued as the post-transparency Fossil-CMB temperature law.

With the corresponding three-dimensional geometric eigenmode correction, the TFT annealing solution reaches the adopted neutral-hydrogen stability and optical-transparency scale near

\begin{equation}
    T_{\mathrm a}
    \sim
    3000\ \mathrm{K},
    \qquad
    t_{\mathrm a}
    \simeq
    3.86 \times 10^5\ \mathrm{yr}.
\end{equation}

Thus, within the TFT annealing closure, the familiar ``\(380{,}000\)-year'' epoch acquires a direct mechanical interpretation: it is the time required for the post-decompression State-I medium and its coupled radiation--matter environment to anneal from the primordial structural state to the transparency regime. The complete derivation is given in \hyperref[app:lattice_annealing]{Appendix~\ref*{app:lattice_annealing}}.

\subsection*{Mechanical Control of the Epoch of Recombination}

This lattice annealing does not replace baryonic recombination; \textbf{it supplies the thermal evolution that governs when recombination becomes possible}. The primordial baryonic plasma---the topological shrapnel produced during Catastrophic Decompression---remains thermally coupled to the surrounding radiation and Active Vacuum during the early annealing phase.

As the coupled system cools toward the few-thousand-kelvin regime, stable neutral hydrogen becomes energetically favored, the free-electron population falls, and the Thomson-scattering opacity decreases rapidly. Once the photon mean free path becomes sufficiently large, the medium becomes optically transparent and the primordial thermal photon population enters its post-transparency Fossil-CMB evolution.

The \(3.86 \times 10^5\)-year result therefore connects the stages of the same TFT mechanical history:

\begin{equation}
\begin{aligned}
    \text{lattice annealing}
    &\longrightarrow
    \text{baryonic recombination}
    \\
    &\longrightarrow
    \text{optical transparency}
    \\
    &\longrightarrow
    \text{Fossil-CMB evolution}.
\end{aligned}
\end{equation}